\section{Discussion}\label{sec:discussion}

\paragraph{Where the steering effect works, it is not vocabulary suppression.}
The deflationary hypothesis that motivated this work is, for the one construction with a real behavioral effect, decisively false. The mass-mean vector's truthfulness gain survives certified excision of its entire first-order vocabulary readout, at every excision rank we tested, under norm matching, across three independent constructions of the token set, and under paraphrase shift. The aligned subspace is not even special: random excisions of equal rank cost the same. The effect is carried by what downstream layers do with the injected direction, not by what the unembedding reads off it directly. This is the generalization-relevant mechanism: a readout-level tilt has no reason to transfer across phrasings, while an upstream signal consumed by later computation does, and $\rho_{\mathrm{OOD}} \approx 1$ on paraphrases is exactly that prediction confirmed.

\paragraph{Where steering is vocabulary-level, it does not work.}
The CAA honesty vector presents the complementary lesson. It moves honesty-coded vocabulary strongly --- over a logit of curated honest-shift at the answer position --- while leaving truthfulness flat to slightly negative on held-out questions. A benchmark suite that scored word choice, persona, or self-reported honesty would crown it; TruthfulQA with a held-out split does not. The two families were built from different supervision: contrast prompts that instruct the model to \emph{behave} honestly versus statements labeled by what is \emph{true}. The construction inherited exactly what its supervision specified --- a register, not an epistemic state. For safety applications this distinction is the entire point: an intervention that makes a model sound more honest without making it more truthful is not a partial success, it is a calibrated failure mode that benchmark-level evaluation will systematically misread as alignment progress.

\paragraph{The instrument must be checked before the mechanism is probed.}
Our naive sweep produced a published-looking result --- a benchmark gain and a clean $\rho \approx 1$ ``deep effect'' --- from a vector that multiplies held-out perplexity seventy-fold. Nothing in the standard evaluation pipeline flags this: teacher-forced multiple choice is structurally insensitive to generation collapse, and a ratio statistic happily divides one artifact by another. The two preconditions we impose (a denominator bounded away from zero at a coherence-gated operating point, and an aligned-versus-random contrast) are cheap, and our results suggest they should be standard practice for any causal claim built on steering interventions. Steering-vector strength sweeps that report only benchmark deltas are, on this evidence, not interpretable.

\paragraph{Implications for readout-based interpretability.}
The lens trajectories show downstream blocks re-synthesizing the suppressed readout within two layers, and the per-prompt word-level depth ratios show that most of each vector's vocabulary effect is manufactured downstream even when the direct path is available. This is a concrete, quantified instance of the non-identifiability that \citet{venkatesh2026nonid} prove in the abstract: the component of a steering vector visible to the unembedding is a poor guide to what the vector does. Methods that interpret steering vectors by reading them through the vocabulary projection --- logit-lens decoding of steering directions, vocabulary-space vector arithmetic --- are measuring the component our experiments show to be causally inert.
